\chapter{Introduction}

\begin{itemize}
\item Language models can state false information with confidence. This is usually called
  \textbf{hallucination} \cite{ji2023survey}.
\item \textbf{Retrieval-augmented generation (RAG)} addresses this by giving the model passages
  from a document collection and asking it to answer from them \cite{lewis2020rag}.
\item RAG does not remove the problem. When the retriever returns the wrong passages, the
  generator builds its answer on bad evidence \cite{yan2024crag}, and the answer can include
  claims that the passages do not back \cite{niu2024ragtruth}.
\item This matters most where people rely on an answer without checking the source, because the
  presence of retrieved passages makes an answer look better founded than it may be.
\item Two published systems, \textbf{Self-RAG} \cite{asai2024selfrag} and
  \textbf{CRAG} \cite{yan2024crag}, decide when to trust retrieval, but both need a fine-tuned
  model. That rules them out when a model is only reachable through an API, or when there is no
  budget for training.
\item Retrieval pipelines already compute a relevance score for every passage, and an answer can
  be compared with the passages it was given. Both could act as a confidence signal, with no
  extra training.
\item This thesis \textbf{measures} how retrieval quality affects hallucination for two small
  open-weight language models, and tests a \textbf{training-free rule}: withhold the answer when
  such a confidence signal is low.
\item \textbf{Structure:} Chapter~\ref{chap:problem} states the problem and the research
  questions, Chapter~\ref{chap:literature} surveys the literature,
  Chapter~\ref{chap:methodology} describes the methodology, Chapter~\ref{chap:results} reports
  the results and Chapter~\ref{chap:conclusion} concludes.
\end{itemize}
